\documentclass[10pt,a4paper]{amsart}
\usepackage[margin=19mm]{geometry}
\usepackage{amsmath,amssymb,mathtools}
\usepackage[scr=rsfs,cal=euler]{mathalfa}
\usepackage{xfrac}
\usepackage[unicode,hidelinks,bookmarks=false]{hyperref}
\newcommand{\ie}{i.e.\ }
\newcommand{\cf}{cf.\ }
\newcommand{\resp}{respectively\ }
\newcommand{\Subsection}{\subsection}
\providecommand{\nobreakdash}{\nobreak\hbox{-}}
\setlength{\emergencystretch}{2em}
\multlinegap0pt
\usepackage[shortlabels]{enumitem}
\usepackage{mathtools}
\usepackage{amssymb}
\usepackage{xcolor}
\newtheorem{proposition}{Proposition}
\newtheorem{theorem}{Theorem}
\providecommand{\Cov}{\operatorname{Cov}}
\providecommand{\E}{\mathbb{E}}
\providecommand{\Leh}{\mathscr{L}}
\providecommand{\R}{\mathbb{R}}
\providecommand{\dd}{\,\mathrm{d}}
\providecommand{\eqdef}{\mathrel{\vcentcolon=}}
\providecommand{\lr}{\mathrm{lr}}
\title{Lehmer--Lambert Distribution}
\author{Masoud Ataei, Vladimir V. Vinogradov}
\date{Source: arXiv:2609.12150v1 (CC BY 4.0)}
\begin{document}
\maketitle

\noindent\emph{Source and licence.} Lehmer--Lambert Distribution. Version arXiv:2609.12150v1 (CC BY 4.0), distributed by arXiv under the Creative Commons Attribution 4.0 International licence, http://creativecommons.org/licenses/by/4.0/, as recorded in the arXiv licence metadata for this exact version and shown on the abstract page https://arxiv.org/abs/2609.12150v1. Full licence notice: \texttt{licenses/arxiv-2609.12150-CC-BY-4.0.txt}.\par\smallskip

\noindent\textit{Editorial scope.} Counted unit: the article's theorem thm:shape (source0.tex lines 719--746). The article's complete original proof of it is delivered with it (source0.tex lines 1464--1534, one \texttt{proof} environment of 71 lines, which the article places away from the statement). No mathematics is written, repaired or re-derived: the statement and the proof are the article's own lines, in the article's own notation, and no retained line is substituted or re-flowed.  Retained uncounted as the source-local context the proof needs: lines 507--510; lines 512--514; lines 590--594; lines 597--656; lines 714--717; lines 918--977.  Nothing was omitted from any retained span, so the package carries no omission record.  No retained line contains a citation, so the wrapper carries no bibliography and no bibliography entry.  No retained line contains a figure environment, an image-inclusion command or drawing code, so no image is delivered and the package declares no asset dependency.  Editorial formatting only, in addition: an AMS \texttt{amsart} preamble that restates the article's own declarations and macros used by the retained lines, namely the environments \texttt{proposition}, \texttt{theorem} on separate counters, together with the standard packages those lines need.  The retained environments print in this document as Proposition 1, Theorem 1 in order of appearance, whereas the article prints the same items with the numbering of its own full document.  The complete native source of the article as distributed in the arXiv e-print for this exact version is bundled unchanged as \texttt{sources/210084/source0.tex}.
			\begin{equation}\label{eq:lgdensity}
				f_Z(z)=C\Bigl(\frac1\alpha+\beta z\Bigr)z^{\frac1\alpha-1}e^{\beta z},
				\qquad 1<z<L_+ .
			\end{equation}

			\begin{equation}\label{eq:ll}
				f_S(s)=f_Z\bigl(\Leh_\nu(s)\bigr)\,\Cov_{\nu_{s-1}}\bigl(X,\log X\bigr).
			\end{equation}

	\begin{equation}\label{eq:gaps}
		r_+\eqdef\tau\log\frac{x_{(n)}}{x_{(n-1)}},
		\qquad
		r_-\eqdef\tau\log\frac{x_{(2)}}{x_{(1)}},
	\end{equation}

	\begin{proposition}[Special cases and limits]
		\label{prop:limits}\leavevmode
		\begin{enumerate}[label=\textup{(\roman*)},leftmargin=2.2em]
			\item At $\beta=0$,
			\begin{equation}\label{eq:powermap}
				F_S(s)=\frac{\Leh_\nu(s)^{1/\alpha}-1}{L_+^{1/\alpha}-1},
			\end{equation}
			and $Z^{1/\alpha}$ is uniform on $(1,L_+^{1/\alpha})$. At
			$\alpha=1$, $\beta=0$ the random Lehmer mean is uniform.
			\item As $\alpha\to\infty$ at $\beta=0$,
			\begin{equation}\label{eq:jeffreysF}
				F_S(s)\longrightarrow\frac{\log\Leh_\nu(s)}{\log L_+} .
			\end{equation}
			This is the \emph{Jeffreys law} of $\nu$, named after the
			Jeffreys divergence that appears in its density in the next
			section and not after the Jeffreys prior of the escort family,
			which is proportional to the square root of the Fisher
			information. Applied to $\mu$ without standardization it is the
			parameter-free law
			\begin{equation}\label{eq:jeffreysraw}
				F_S(s)=\frac{\Lambda_\mu(s)-\log m}{\ell_\mu} .
			\end{equation}
			Under it $\log\Leh_\mu(S)$ is uniform on $(\log m,\log M)$ and
			\begin{equation}\label{eq:logmean}
				\E\bigl[\Leh_\mu(S)\bigr]=\frac{M-m}{\ell_\mu},
			\end{equation}
			the logarithmic mean of the extremes.
			\item Let $L_+=1+\Omega$ and $\Leh_\rho(y)=ye^{y}$, the reference curve
			with $\sigma^{2}=1$. Then
			\begin{equation}\label{eq:canonical}
				F_S(s)=\Leh_\rho\bigl(\Leh_\nu(s)-1\bigr)
			\end{equation}
			is a distribution function, the \emph{canonical Omega law}, and
			$Y=\Leh_\nu(S)-1=W_0(U)$ with $U$ uniform on $(0,1)$.
			\item Let $\mu$ be finitely supported, let $w_M,w_M'$ be the masses
			of the largest support point and of the largest support point
			below it, and let $r_+$ be the upper gap of definition
			\eqref{eq:gaps}. Put
			\begin{equation}\label{eq:lambdaa}
				\lambda=\beta+\frac{1}{\alpha L_+},
				\qquad
				a_+=L_+\bigl(e^{r_+}-1\bigr)\frac{w_M'}{w_M} .
			\end{equation}
			As $\lambda\to\infty$,
			\begin{equation}\label{eq:gumbel}
				\lambda(L_+-Z)\xrightarrow{d}\mathrm{Exp}(1),
			\end{equation}
			\begin{equation}\label{eq:gumbel2}
				r_+S-\log(a_+\lambda)\xrightarrow{d}\mathrm{Gumbel}.
			\end{equation}
			This covers $\beta\to\infty$ and $\alpha\downarrow0$ alike.
			\item As $L_+\downarrow1$ at fixed $(\alpha,\beta)$, $\tau S$
			converges in law to the variable $T$ with density
			\begin{equation}\label{eq:infoprofile}
				f_T(u)=\frac{K_\mu''(u)}{\ell_\mu} .
			\end{equation}
			This is the Fisher information of the escort family of the raw
			signal, read as a probability density.
		\end{enumerate}
	\end{proposition}

	\begin{equation}\label{eq:psi}
		\psi(z)\eqdef\frac{\dd}{\dd z}\log f_Z(z)
		=\frac{\alpha\beta}{1+\alpha\beta z}+\frac{1-\alpha}{\alpha z}+\beta .
	\end{equation}

	\begin{theorem}[Jeffreys law]
		\label{thm:jeffreys}
		Let $S$ have the Jeffreys law of $\mu$, with distribution
		function \eqref{eq:jeffreysraw}.
		\begin{enumerate}[label=\textup{(\roman*)},leftmargin=2.2em]
			\item The density is a normalized Jeffreys divergence between
			consecutive escorts,
			\begin{equation}\label{eq:jeffreysdensity}
				f_S(s)=\frac{J(\mu_{s-1},\mu_s)}{\ell_\mu}
				=\frac1{\ell_\mu}\int_{s-1}^{s}I(u)\dd u,
				\qquad
				\int_\R I(u)\dd u=\ell_\mu .
			\end{equation}
			\item $S=T+U$ with $T$ and $U$ independent, $U$ uniform on $(0,1)$
			and $T$ the \emph{information profile} of the signal,
			\begin{equation}\label{eq:Tlaw}
				f_T(u)=\frac{I(u)}{\ell_\mu},
			\end{equation}
			\begin{equation}\label{eq:Tlaw2}
				F_T(u)=\frac{K_\mu'(u)-\log m}{\ell_\mu} .
			\end{equation}
			Hence
			\begin{equation}\label{eq:unitincrement}
				f_S(s)=F_T(s)-F_T(s-1)\le1,
			\end{equation}
			and the stationary points of the density of $S$, among them its
			modes, are the points where $I(s)=I(s-1)$. The
			cumulants satisfy $\kappa_j(S)=\kappa_j(T)+B_j/j$, with $B_j$ the
			Bernoulli numbers and $B_1=\tfrac12$.
			\item If $\mu$ has atoms at both extremes with probabilities
			$p_{\min},p_{\max}$, then
			\begin{equation}\label{eq:surprisal}
				\E[T^{+}]=-\frac{\log p_{\max}}{\ell_\mu},
				\qquad
				\E[T^{-}]=-\frac{\log p_{\min}}{\ell_\mu},
			\end{equation}
			\begin{equation}\label{eq:surprisal2}
				\E[S]=\frac12+\frac{\log(p_{\min}/p_{\max})}{\ell_\mu} .
			\end{equation}
			For a sample of $n$ observations with distinct extremes,
			\begin{equation}\label{eq:samplemean}
				\E[S]=\frac12,
				\qquad
				\E|T|=\frac{2\log n}{\ell_\mu},
			\end{equation}
			whatever the interior of the sample. For $\mu=\delta_m+\delta_M$, $T$
			is logistic with scale $1/\ell_\mu$.
			
			\item If $\mu$ is finitely supported, the right tail is exponential
			with rate $\log(x_{(n)}/x_{(n-1)})$. If
			instead $\mu$ has a density $f(x)\sim c(M-x)^{k}$ as $x\uparrow M$,
			then
			\begin{equation}\label{eq:cauchytail}
				1-F_S(s)\sim\frac{k+1}{\ell_\mu s},
			\end{equation}
			and $\E[S^{+}]=\infty$. The mirror statements hold at $m$. For
			any member of the family the same dichotomy holds with the
			constant $f_Z(L_+)L_+(k+1)$ in place of $(k+1)/\ell_\mu$.
		\end{enumerate}
	\end{theorem}

	\begin{theorem}[Shape]\label{thm:shape}\leavevmode
		\begin{enumerate}[label=\textup{(\roman*)},leftmargin=2.2em]
			\item For $\alpha\le1$ the density $f_Z$ is nondecreasing and
			log-concave and $G$ is convex, whatever $\beta\ge0$ and $L_+$,
			strictly so unless $\alpha=1$ and $\beta=0$. For $\alpha>1$ this
			fails when $\beta$ is small.
			\item For $\alpha\le1$ every interior mode of $S$ satisfies
			\begin{equation}\label{eq:modeeq}
				\frac{\Leh_\nu''(s)}{\Leh_\nu'(s)}=-\Leh_\nu'(s)\,\psi\bigl(\Leh_\nu(s)\bigr)\le0 ,
			\end{equation}
			with equality only when $\alpha=1$ and $\beta=0$, so the modes
			lie in the closed concavity episodes of the curve.
			\item The family increases in the likelihood-ratio order as $\beta$
			increases, and as $\alpha$ decreases within the range
			$\alpha\le1$. Every member dominates the Jeffreys law $S_J$ of
			$\nu$,
			\begin{equation}\label{eq:lrjeffreys}
				S\ge_{\lr}S_J .
			\end{equation}
			For a sample with distinct extremes, $\E[S]\ge\tfrac12$.
			\item For two standardized measures $\nu,\nu'$ at the same level,
			\begin{equation}\label{eq:kolmogorov}
				\sup_s\bigl|F_S(s)-F_{S'}(s)\bigr|
				\le\Bigl(\sup_{[1,L_+]}f_Z\Bigr)\sup_s\bigl|\Leh_\nu(s)-\Leh_{\nu'}(s)\bigr| ,
			\end{equation}
			and $\sup f_Z=CG'(L_+)$ when $\alpha\le1$.
		\end{enumerate}
	\end{theorem}

		\begin{proof}[Proof of Theorem~\ref{thm:shape} (Shape)]\leavevmode
			
			\noindent(i) From the density \eqref{eq:lgdensity},
			\begin{equation}\label{eq:psiprime}
				\psi'(z)=-\frac{(\alpha\beta)^{2}}{(1+\alpha\beta z)^{2}}-\frac{1-\alpha}{\alpha z^{2}} .
			\end{equation}
			For $\alpha\le1$ the three terms of $\psi$ in formula \eqref{eq:psi}
			are nonnegative and the two terms of $\psi'$ are nonpositive,
			whatever $\beta\ge0$, so $f_Z$ is nondecreasing and log-concave;
			both are strict unless $\alpha=1$ and $\beta=0$, when $\psi\equiv0$.
			For $\alpha>1$ the middle term of $\psi$ equals $1/\alpha-1<0$ at
			$z=1$ and dominates when $\beta$ is small, while for large $\beta$
			the last term dominates. The hazard properties
			are the standard ones of log-concave densities on an interval, the
			divergence following from $A-G(z)\to0$. Convexity of $G$ follows
			from
			\begin{equation}\label{eq:Gsecond}
				G''=\bigl[a(a-1)z^{-2}+2a\beta z^{-1}+\beta^{2}\bigr]G,
			\end{equation}
			which is nonnegative for all $z\ge1$ and all $\beta\ge0$ if and only
			if $a\ge1$, since at $\beta=0$ it has the sign of $a-1$.
			
			\noindent(ii)
			Logarithmic differentiation of the density \eqref{eq:ll} gives
			\begin{equation}\label{eq:logfS}
				(\log f_S)'=\psi(\Leh_\nu)\Leh_\nu'+\frac{\Leh_\nu''}{\Leh_\nu'} .
			\end{equation}
			At a stationary point $\psi(\Leh_\nu)\Leh_\nu'\ge0$, since
			$\psi\ge0$ for $\alpha\le1$ and $\Leh_\nu'>0$, so $\Leh_\nu''\le0$.
			The inequality is strict unless $\psi\equiv0$, that is unless
			$\alpha=1$ and $\beta=0$, when the density of $S$ is proportional
			to $\Leh_\nu'$ and its modes are the inflection points at which
			the curve turns concave.
			
			\noindent(iii) For
			$\beta'>\beta$ the ratio $f_Z(z;\alpha,\beta')/f_Z(z;\alpha,\beta)$
			is proportional to
			\begin{equation}\label{eq:lrbeta}
				\frac{1+\alpha\beta'z}{1+\alpha\beta z}\,e^{(\beta'-\beta)z},
			\end{equation}
			a product of increasing functions. For $\alpha'<\alpha$ the log-derivative in $z$ of $f_Z(z;\alpha',\beta)/f_Z(z;\alpha,\beta)$
			equals
			\begin{equation}\label{eq:lralpha}
				\frac{\alpha-\alpha'}{\alpha\alpha'z}\Bigl[1-\frac{\alpha\alpha'\beta z}{(1+\alpha'\beta z)(1+\alpha\beta z)}\Bigr],
			\end{equation}
			whose bracket is positive when $\alpha'\le1$, because
			\begin{equation}\label{eq:lralphabracket}
				(1+\alpha'\beta z)(1+\alpha\beta z)-\alpha\alpha'\beta z
				=1+\alpha'\beta z+\alpha\beta z(1-\alpha')+\alpha\alpha'\beta^{2}z^{2}>0 .
			\end{equation}
			When $\alpha'>1$ the third term is negative and the left side of
			identity \eqref{eq:lralphabracket} is negative on an interval of
			values of $\beta z$ as soon as
			$(\alpha+\alpha'-\alpha\alpha')^{2}>4\alpha\alpha'$, so the
			restriction cannot be dropped. Both laws
			of $S$ are images of laws on $(1,L_+)$ under the common bijection
			$\Leh_\nu^{-1}$. So the likelihood-ratio order transfers to the
			suddency line, and it implies the usual stochastic order, hence
			the monotonicity of quantiles and means. The Jeffreys law of $\nu$
			makes $\log Z$ uniform on $(0,\ell)$ by
			Proposition~\ref{prop:limits}(ii). So the likelihood ratio of $S$
			with respect to $S_J$ is $\ell zf_Z(z)$ at $z=\Leh_\nu(s)$, which
			is $\ell C(a+\beta z)z^{a}e^{\beta z}$, increasing in $z$ and hence
			in $s$. Stochastic order gives $\E[S]\ge\E[S_J]=\tfrac12$ by
			Theorem~\ref{thm:jeffreys}(iii) applied to $\nu$.
			
			\noindent(iv) is the mean value theorem applied to the map
			$z\mapsto C(G(z)-e^{\beta})$, whose derivative is $f_Z$; for
			$\alpha\le1$ the density is nondecreasing by (i), so its supremum is
			$f_Z(L_+)=CG'(L_+)$.
		\end{proof}

\end{document}
